% ===================================================================
\section{The Spectral Attractor}
\label{sec:attractor}

\subsection{Universal Convergence}

Paper~3 reported that spectral-only training (contrastive loss
disabled) at $n \in \{3, 4, 8\}$ converges to Bridge Fiedler
${\sim}0.09$. We extend the ablation to $n = 12$ and report the
complete results in Table~\ref{tab:attractor}.

\begin{table}[H]
\centering
\caption{Spectral attractor: Bridge Fiedler eigenvalue under
  spectral-only training (no contrastive loss) across channel
  counts. All experiments use TinyLlama-1.1B, identity
  initialization, and identical hyperparameters. H-ch6 (H3) is
  the $n = 6$ experiment with contrastive loss active, shown for
  bifurcation comparison.}
\label{tab:attractor}
\small
\begin{tabular}{lcccccl}
\toprule
ID & $n$ & Contrastive & Steps & Fiedler & $\rho$ & BD \\
\midrule
H-ch3  & 3  & No  & 10{,}000 & 0.0951  & N/A       & No \\
H-ch4  & 4  & No  & 10{,}000 & 0.0836  & ${\sim}1$:1 & No \\
H-ch8  & 8  & No  & 10{,}000 & 0.0944  & ${\sim}1$:1 & No \\
H-ch12 & 12 & No  & 10{,}000 & 0.1019  & ${\sim}1$:1 & No \\
\midrule
H3 (H-ch6) & 6 & \textbf{Yes} & 10{,}000 & \textbf{0.00009} & \textbf{70{,}404:1} & \textbf{Yes} \\
\bottomrule
\end{tabular}
\end{table}

The spectral-only experiments converge to a Fiedler band of
$0.0836$--$0.1019$ (Figure~\ref{fig:spectral_attractor})---a 10.4\%
coefficient of variation across a $4\times$ span of channel counts
($n = 3$ to $n = 12$). The
attractor value clusters near the spectral loss target of
$\lambda_2^* = 0.1$ described in Paper~3. This convergence is
\emph{channel-count-independent}: the bridge reaches the same
connectivity equilibrium whether it has 3, 4, 8, or 12~channels.

At the attractor, bridge structure is \emph{connected but
undirected}. Co-planar/cross-planar ratios are ${\sim}1$:1 (where
applicable), and no block-diagonal structure emerges. The spectral
loss distributes connectivity uniformly across all channel pairs;
without contrastive specification, there is no reason for the bridge
to favor any particular pair over another.

H-ch12 ($n = 12$) is the most demanding case. With
$\binom{12}{2} = 66$ off-diagonal pairs sharing a fixed connectivity
budget, convergence is slower but the bridge reaches Fiedler $0.1019$
at step~10{,}000---confirming the attractor holds at $4\times$ the
minimum channel count. The slight upward trend (0.1019 vs.\
$0.0836$--$0.0951$ for smaller~$n$) suggests a weak positive
relationship between channel count and final Fiedler, but the effect
is small relative to the attractor's stability across the full range.

\subsection{Interpretation}

The spectral attractor has a straightforward interpretation. The
spectral loss penalizes deviations of the bridge's Fiedler eigenvalue
from a target~$\lambda_2^*$. Under spectral regularization alone, the
bridge distributes coupling uniformly across all off-diagonal entries:
this is the minimum-energy configuration that achieves the target
$\lambda_2$. Adding more channels does not change this equilibrium
because each new channel adds both off-diagonal entries (coupling
budget to distribute) and the same spectral target (connectivity
to achieve). The ratio of budget to target remains constant.

The attractor is the bridge's \emph{unprogrammed state}---the
resting connectivity it achieves before any topological direction
is specified. It is analogous to a uniform probability distribution:
maximum entropy given the connectivity constraint. The spectral
loss provides the constraint (how much total connectivity); the
contrastive loss provides the structure (where connectivity
concentrates).

\begin{figure}[t]
  \centering
  \includegraphics[width=\textwidth]{figures/fig_spectral_attractor.pdf}
  \caption{Spectral attractor universality across channel counts.
    Fiedler eigenvalue converges to ${\approx}\,0.09$--$0.10$ under
    spectral-only training for $n \in \{3, 4, 8, 12\}$---a 10.4\%
    band across a $4\times$ range of channel counts. The attractor is
    the bridge's natural connectivity equilibrium in the absence of
    contrastive topology specification.}
  \label{fig:spectral_attractor}
\end{figure}

\subsection{The Bifurcation}
\label{sec:bifurcation}

The spectral attractor provides the baseline against which
contrastive topology programming is measured. Adding contrastive
loss breaks the attractor by specifying directional coupling
preferences, producing a sharp bifurcation in spectral connectivity.

The bifurcation is extreme at both confirmed dimensionalities:

\begin{itemize}
  \item \textbf{RD ($n = 6$):} Fiedler drops from $0.0918$
    (spectral-only, H-ch4 at $n = 4$, nearest comparable) to
    $0.00009$ (contrastive active). A factor of $1{,}020\times$.
    Co/cross: ${\sim}1$:1 $\to$ $70{,}404$:1.
  \item \textbf{Tesseract ($n = 8$):} Fiedler drops from $0.0944$
    (spectral-only, H-ch8) to $0.000191$ (contrastive active,
    T-001r2). A factor of ${\sim}470\times$. Co/cross: ${\sim}1$:1
    $\to$ $41{,}564$:1.
  \item \textbf{Block-diagonal classification:} 0\% (spectral-only
    at all channel counts) $\to$ 100\% (contrastive active at
    $n = 6$ and $n = 8$).
\end{itemize}

This bifurcation decomposes the Steersman's effect into two
orthogonal control channels:

\begin{enumerate}
  \item \textbf{Spectral loss $\to$ connectivity.} The spectral
    component drives the bridge toward a target algebraic
    connectivity ($\lambda_2 \approx 0.1$). It makes the bridge
    \emph{connected}---channels interact---but imposes no directional
    preference. The result is the spectral attractor: a connected,
    isotropic bridge.
  \item \textbf{Contrastive loss $\to$ topology.} The contrastive
    component specifies which connections to strengthen (co-axial
    pairs) and which to suppress (cross-axial pairs). It transforms
    isotropic connectivity into directed topology. The result is
    block-diagonal structure whose block count matches the number
    of co-axial pairs.
\end{enumerate}

Neither component alone produces block-diagonal structure. Spectral
loss alone produces the attractor. Contrastive loss alone (without
spectral regularization) was not tested in this series but is a
natural follow-up. The two components together produce a controlled
bifurcation: connectivity plus direction yields topology.

\subsection{Comparison: Bifurcation Across Dimensionalities}

Table~\ref{tab:bifurcation} summarizes the confirmed bifurcation
events. The bifurcation ratio is smaller at $n = 8$ than at $n = 6$,
likely because the larger number of cross-axial entries at $n = 8$
(24~vs.~12) creates more total suppression pressure, producing a
slightly higher residual Fiedler at convergence. Both bifurcations
are of the same order of magnitude ($10^{2\text{--}3}\times$) and
both produce complete block-diagonal structure (100\% classification).

\begin{table}[H]
\centering
\caption{Bifurcation comparison: spectral attractor broken by
  contrastive topology specification at two dimensionalities.}
\label{tab:bifurcation}
\small
\begin{tabular}{lccccc}
\toprule
Geometry & $n$ & Spectral Fiedler & Contrastive Fiedler & Bifurcation & $\rho$ \\
\midrule
RD (3D)       & 6 & 0.0918 & 0.00009  & $1{,}020\times$ & 70{,}404:1 \\
Tesseract (4D) & 8 & 0.0944 & 0.000191 & ${\sim}470\times$  & 41{,}564:1 \\
\bottomrule
\end{tabular}
\end{table}
